\section{Availability}
\label{sec:availability}
\index{Availibility}
Not all equipment is always available everywhere.
If you want to test whether an item is accessible,
	you can determine by rolling as follows.
\par
\begin{enumerate}
	\vspace{-10mm}
	\setlength\itemsep{-10mm}
	\item Roll a D5.
	\item Subtract a number based on where you’re at.
	\begin{sitemize}
		\item +1 for any item in an unspecialized shop (like a pawn shop).
		\item +3 for a misplaced item in a specialized shop, e.g. chrome from a gun fixer.
		\item Buying from a megacorp is unmodified.
		\item Buying from a nomad trader or a bazaar infers +1.
		\item Buying from a (proverbial or literal) crack den infers +2.
	\end{sitemize}
	\item If that number \emph{is lower or equal} to the item's Availability score, it is available.
	Otherwise it is not.
\end{enumerate}
\par
Such a roll is meant to be made when quickly looking through a merchant's inventory.
Given enough time any specific item can probably be acquired.
\\%
Some objects may also simply be unavailable at times.
In such a case no roll is even warranted.
An example are implants with the \emph{Perfect Fit} modification.